\clearpage
\section{Appendix: Liturgical Resolution}
\label{sec:liturgical-resolution}

{\small
\textbf{Celebration.} The Twenty-ninth Sunday in Ordinary Time, a Sunday in
Ordinary Time at place~6 of the table of liturgical days; green; Gloria and
Creed. The formula is \bid{PC-S55-A}, the Year~A target of the parent
\bid{PC-S55}; it has the single Sunday form. A Mass anticipated on Saturday
evening uses the same formulary and readings, and there is no Vigil form.

\textbf{Books, calendar and territory.} The \work{Roman Missal}, Third
Edition, for use in the dioceses of the United States of America (2011), whose
Latin base is the \work{Missale Romanum}, editio typica tertia (2002), as
reprinted with emendations in 2008; the \work{Lectionary for Mass for Use in
the Dioceses of the United States of America}, second typical edition,
volume~I, no.~145, which follows the \work{Ordo lectionum Missae}, editio
typica altera (1981); the \work{General Instruction of the Roman Missal} in the
United States edition of 2011 with the emendations of 2021; the
\work{Universal Norms on the Liturgical Year and the Calendar} and the General
Roman Calendar, in the Latin printed in the 2002 Missal; and the
\work{Liturgical Calendar for the Dioceses of the United States of America} for
2026. Only the national calendar is applied. No diocesan, religious, parish,
titular, dedication or patronal calendar was supplied, and none is assumed.

\textbf{Cycle and occurrence.} Year~A of the Sunday cycle runs from
30~November 2025 to 22~November 2026, and the national calendar enters Sunday
18~October 2026 as the Twenty-ninth Sunday in Ordinary Time, green, with the
readings of Lectionary no.~145 and Psalter Week~I. The General Roman Calendar
inscribes the Feast of Saint Luke, Evangelist, on 18~October. A feast of a
saint of the General Calendar ranks below a Sunday in Ordinary Time, and under
\work{Universal Norms} 60 it is omitted in this year: it is neither
transferred nor commemorated, and it supplies no text to this Mass. The weekday
cycle governs no Sunday.

\textbf{The Missal formulary and the three cycles.} The antiphons and orations
belong to the Twenty-ninth Week in Ordinary Time and are common to Years A, B
and~C; their text is recorded once, in the collection's record for that week.
The readings, psalm and Alleluia verse are those of Year~A alone. The second
reading begins the Lectionary's course through the First Letter to the
Thessalonians in Year~A, and the Gospel belongs to its semi-continuous reading
of Matthew.

\textbf{Branches.} No selection among the following was supplied, and none is
inferred. Both Communion antiphons are studied in full; a claim resting on one
of them holds only where it is chosen, and a claim resting on the wording of an
antiphon does not hold where another approved chant replaces it. The last three
branches would replace this formulary entirely, and nothing in this study holds
for a celebration that uses one of them.
\par}

\begin{branchtable}
\bid{communion-antiphon-psalm} & The Missal prints cf.\ Ps 33 (32):18--19 first
& appointed alternative & Communion Antiphon & Unselected; studied in full.\\
\bid{communion-antiphon-gospel} & The Missal prints Mk 10:45 second, under
\latin{Vel} & appointed alternative & Communion Antiphon & Unselected; studied
in full. Rubric~6 for Ordinary Time prefers an antiphon that agrees with the
Gospel of the Mass; this one is drawn from a Gospel, but not from this
Sunday's, and the books do not say whether either antiphon agrees with it.\\
\bid{entrance-or-communion-chant-substitution} & \work{General Instruction} 48
and 87: the Missal or \work{Graduale Romanum} antiphon, the \work{Graduale
Simplex}, a chant from another collection of psalms and antiphons approved by
the Conference of Bishops or the Diocesan Bishop, or another suitable
liturgical chant likewise approved & permitted & Entrance; Communion &
Unselected; the Missal antiphons are the texts studied.\\
\bid{psalm-or-acclamation-chant-substitution} & \work{General Instruction} 61
and 62: a common response and psalm for the time of year when the psalm is
sung; the \work{Graduale} forms or an antiphon and psalm from another approved
collection, but never a song or hymn in place of the psalm & permitted &
Responsorial Psalm; Alleluia verse & Unselected; the Lectionary's psalm and
verse are the texts studied.\\
\bid{offertory-chant} & \work{General Instruction} 74 & permitted &
Preparation of the gifts & Unselected; the Missal prints no Offertory
antiphon.\\
\bid{sprinkling-rite} & \work{General Instruction} 51 and the Order of Mass:
on Sundays the blessing and sprinkling of water may replace the Penitential
Act & ritual substitution & Introductory Rites & Unselected; it brings its own
texts.\\
\bid{compatible-preface-and-eucharistic-prayer} & Rubric~5 for Ordinary Time;
\work{General Instruction} 364--365 & permitted & Eucharistic Prayer &
Unselected; one of the eight Prefaces of the Sundays in Ordinary Time with a
compatible Eucharistic Prayer, as one coupled choice.\\
\bid{eucharistic-prayer-iv-with-preface} & \work{General Instruction} 365~d &
conditional & Eucharistic Prayer & Unselected; Eucharistic Prayer~IV keeps its
own invariable Preface.\\
\bid{concluding-blessing-form} & The Missal's solemn blessings and prayers over
the people, at the priest's discretion & permitted & Concluding Rites &
Unselected; none is appointed to this Sunday.\\
\bid{mass-for-evangelization-of-peoples} & The Missal's Mass for the
Evangelization of Peoples, admitted on Sundays in Ordinary Time when particular
celebrations for the missions are held & conditional & Whole Mass & Not
supplied. The national calendar prints no mission observance for this date,
and no dated witness establishes one; where this Mass is celebrated, the study
does not hold for it.\\
\bid{ritual-or-various-needs-mass} & \work{General Instruction} 372 and 374 &
conditional & Whole Mass & Not supplied; a replacement would need its own
resolution.\\
\bid{local-proper-calendar} & \work{Universal Norms} 59--60 with a particular
calendar; in a church of which Saint Luke is the titular, the Solemnity of the
Title outranks the Sunday & conditional & Whole Mass & Not supplied; a local
celebration would need its own resolution.\\
\end{branchtable}
